\documentclass[10pt]{article}
\usepackage[margin=20mm]{geometry}
\usepackage{amsmath,amssymb,fontspec,unicode-math,xurl,hyperref}
\IfFontExistsTF{DejaVu Serif}{\setmainfont{DejaVu Serif}}{\setmainfont{Latin Modern Roman}}
\setmathfont{latinmodern-math.otf}
\setlength{\parindent}{0pt}\setlength{\parskip}{6pt}\sloppy
\title{Technical Verification of PaperBench Scores and Grading}
\author{AI-assisted technical review}
\date{}
\begin{document}
\maketitle
Abstract: The inspected hierarchical scoring design is coherent, and the count of 8,316 rubric leaves reconciles with the appendix. Two reported means do not reconcile with their displayed runs: Claude BasicAgent’s full table averages to 19.8167\%, rather than the headline 21.0\%, and a Gemini row has an incompatible mean. A small-sample standard-error convention is also unclear. Static inspection of the later released source identifies an automatic-zero path for results available only in the reproduction log and an aggregation definition differing from the paper’s JudgeEval description. These support bounded numerical and implementation concerns, assessed Minor, with historical impact unresolved. They do not establish reversal of model ordering or the human comparison. Adjudication remains pending.\par
\section*{Identity, versions, and review boundaries}
Giulio Starace, Oliver Jaffe, Dane Sherburn, James Aung, Chan Jun Shern, Leon Maksin, Rachel Dias, Evan Mays, Benjamin Kinsella, Wyatt Thompson, Johannes Heidecke, Amelia Glaese, and Tejal Patwardhan, PaperBench: Evaluating AI’s Ability to Replicate AI Research, ICML 2025, PMLR 267 [paper]. Assignment 0217 / I25-43586; OpenReview xF5PuTLPbn. The frozen publisher PDF retrieved 25 October 2026 has 31 pages; appendices begin on p. 13. Locations below use one-based PDF pages, equal to printed local page numbers.\par
The available source is the migrated repository `openai/frontier-evals`, commit `b78dab1da774a6456b1ed9c48c9ff194ee3ff36d`, dated 12 October 2026 [code]. This is later than the paper, not a verified publication-time implementation. The selective lossless snapshot preserves 212 source/common files with archive/member/Git-blob mappings. The preparation manifest records 518 benchmark or rubric dataset members as excluded. Those payloads, agent/human run artifacts and JudgeEval gold-label/result payloads were unavailable for this review. Their absence here is not a claim of nonrelease.\par
The review inspected the original, inspected relevant code statically, independently transcribed and recalculated decisive table entries, and viewed source pages 3, 21, 24 and 25. No author code, agent, LLM judge or API was run; no software or dataset was downloaded. The optional OpenReview endpoint was blocked by HTTP 403. Later revisions and external errata were not searched.\par
\section*{Scoring definition and checks that agree}
A rubric is a tree with nonnegative sibling-relative importance weights. Leaves receive binary scores. For an internal node \(v\) whose children have positive total weight, \[s_v=\frac{\sum_c w_c s_c}{\sum_c w_c}.\] Recursion gives root score \(S=\sum_\ell a_\ell y_\ell\), where \(a_\ell\) is the product of normalized edge weights along the leaf’s path and \(\sum_\ell a_\ell=1\). Induction proves this: a parent forms a convex combination of its children’s convex combinations. Thus \(S\) measures weighted partial completion, not the fraction of papers fully replicated or an unweighted count of independent experiments. Static inspection of the recursive scorer agrees.\par
The main metric is the mean root score across 20 papers. With three complete repeats it is the mean of all 60 scores, regardless of how repeat columns are paired across papers. Independent summation of Table 9, p. 21, gives \[3674\ {\rm Code\ Development}
 +4076\ {\rm Execution}
 +566\ {\rm Result\ Match}=8316.\] The abstract’s count agrees. These leaves remain nested within 20 papers; they are not automatically thousands of independent research tasks.\par
The paper explicitly awards partial credit for source that appears correct, execution, and matching results. Fresh execution uses a separate environment and `reproduce.sh` as entrypoint. Code-Dev deliberately omits execution and is described as a noisy, less robust proxy (p. 4). No claim that Code-Dev proves empirical replication is found.\par
\section*{Supported discrepancies}
Problem 1 (Reported means are internally inconsistent).  The Claude BasicAgent headline and a Gemini row do not follow from the displayed runs.\par
\subsection*{Claude BasicAgent: complete arithmetic and rounding}
The abstract, Section 5.2 and Table 3, p. 6, report \(21.0\%\) for Claude 3.5 Sonnet BasicAgent. Table 15, p. 24, lists 20 papers and three runs for that setup, with no missing entries. Parent transcription of all 60 values, checked against the source image, sums to \(11890\) in units of \(0.001\). Hence \[\bar S=\frac{11890}{1000\cdot60}
       =\frac{1189}{6000}
       =0.1981666667,\qquad
 100\bar S=19.81666667\%.\] Column means are \(18.165\%,21.715\%,19.570\%\). The overall mean is invariant to separate within-paper column reordering.\par
If each score was rounded to three decimals, its error is at most \(0.0005\), so the mean error is at most \(0.05\) percentage points. The discrepancy is \(1.1833333\) percentage points, far outside this bound even allowing one-decimal headline rounding.\par
\subsection*{Gemini: a row contradicted by its own runs}
Table 17, p. 25, lists the `fre` scores as \(0^*,0^*,0.007\) but gives mean \(0.061\) and standard error \(0.009\). Asterisks indicate disqualified scores set to zero under Section 2.5. Their mean is \[(0+0+0.007)/3=0.0023333333,\] not \(0.061\). Excluding both disqualifications would give \(0.007\), still not the displayed mean. A stale pre-disqualification summary is possible, but not established.\par
\subsection*{Impact and conditional correction}
These are Minor numerical inconsistencies. At least one entry in each conflicting set is wrong or belongs to another run collection. Versioned run identifiers and grade records are needed to recover historical truth. If the appendix scores are authoritative, Claude’s mean rounds to \(19.8\%\) and the Gemini row to \(0.002\). These are conditional arithmetic corrections.\par
The recalculated Claude value remains above o1 BasicAgent’s reported \(13.2\%\), so the main-setup ordering does not reverse. No demonstrated effect on the human comparison or benchmark design follows. Major impact is not inferred solely from the headline location.\par
Problem 2 (Uncertainty reporting needs an explicit finite-sample convention).  Per-paper errors differ from customary sample SEM; several aggregate errors require unavailable historical details.\par
For \(n\) independent repeated scores the conventional sample-based SEM is \[{\rm SE}_1=\sqrt{\frac{\sum_j(s_j-\bar s)^2}{n(n-1)}}.\] Using population variance instead yields \[{\rm SE}_0=\sqrt{\frac{\sum_j(s_j-\bar s)^2}{n^2}}
           =\sqrt{\frac{n-1}{n}}\,{\rm SE}_1.\] For three repeats the factor is \(\sqrt{2/3}=0.8164966\).\par
The Table 10 adaptive-pruning row, p. 21, has scores \(0,0.027,0.193\), giving \[{\rm SE}_1=0.0603388579,\qquad
 {\rm SE}_0=0.0492664712.\] The reported \(0.049\) matches the latter. Keeping the disqualified zero exact and allowing the other scores their full rounding intervals still gives sample SEM between \(0.0601099\) and \(0.0605684\).\par
This is a Minor uncertainty-convention concern, not an assertion that every population-variance plug-in estimate is forbidden. The convention should be declared. If customary sample SEM is intended, use denominator \(n(n-1)\). The later `paperbench/metrics.py:80--148` uses `ddof=1` for complete-run and per-paper statistics; this does not identify the historical implementation.\par
Aggregate uncertainty also depends on the replication unit. Treating Table 15 columns as complete benchmark repeats gives sample SEM \(1.0322\) percentage points. Treating papers as fixed strata with independent within-paper repeats gives \[\frac1{20}\sqrt{\sum_{p=1}^{20}\frac{s_p^2}{3}}
 =0.8505\ {\rm percentage\ points}.\] The paper reports \(0.8\). Rounding and undisplayed raw precision could affect this close stratified comparison, so it is not used as a decisive contradiction. Other aggregate errors remain unresolved: for example, o1 BasicAgent’s displayed complete-run columns give \(1.022\) percentage points, versus \(0.3\) in Table 3.\par
Historical run grouping, revised grades and the exact estimator are needed before deciding which aggregate uncertainty is correct. This is Needs Evidence, not an established significance failure. No confidence interval or ranking reversal is inferred.\par
Problem 3 (The released grader rejects log-only result evidence).  The pinned 2026 source automatically zeros Result Analysis leaves when genuine results appear only in `reproduce.log`, although the paper permits that evidence channel.\par
Section 2.4, p. 3, and Table 1, p. 4, allow Result Match grading using `reproduce.sh`, `reproduce.log` and reproduction outputs. Appendix D, p. 15, describes timestamp filtering for additional result files, without requiring every result to be saved separately.\par
The current implementation calls this category `Result Analysis`. Relative to `project/paperbench/paperbench/`, the control flow is:\par
1. `reproduce.py:47--56` creates a timestamp file and pipes runtime output into `reproduce.log`.\par
2. `judge/base.py:70--101` reads the script and log separately and establishes the reproduction timestamp.\par
3. `judge/simple.py:247--336` filters document/table files. Allowed extensions exclude `.log` and `.creation\_time`; the script and log are explicitly handled separately.\par
4. If all retained files predate reproduction, including when the list is empty, lines 328–334 set `reproduce\_touched\_files=False`.\par
5. Lines 557–568 return valid score zero for every Result Analysis leaf without asking the judge to inspect the log.\par
A complete witness is a submission with preexisting `reproduce.sh` and `method.py`, optionally an old README, that genuinely recomputes results and prints the full numerical evidence to standard output. It creates no new file with an eligible extension; temporary computation may be in memory or outside the submission. The wrapper creates a fresh log and timestamp file, but neither enters the eligible list. With no README the list is empty; with an old README all retained times predate reproduction. Every Result Analysis leaf is therefore zeroed, even for a criterion fully evidenced in the fresh log. Metadata in `reproduce.py:225--241` is written outside the submission and does not rescue the branch.\par
This is a static control-flow proof, not author-program execution or an API test. It establishes a Minor version-specific protocol discrepancy and a deterministic false-negative mode. It does not establish that a historical run had this form, how much rubric weight was affected, or that the reported F1 is false. An imperfect automatic judge need not be correct on every example to achieve its reported average accuracy.\par
A bounded correction is to remove the implication that no eligible auxiliary file means no reproduced result and allow fresh log evidence to be assessed. A log must not grant automatic credit; correctness and provenance still require judging. This removes the forced-zero condition only. No source was changed, and no complete accuracy or adversarial-robustness repair is certified.\par
\section*{Judge validation and comparison scope}
\subsection*{The current JudgeEval aggregation is a different statistic}
Section 4.2, p. 6, describes macro-averaging across papers. The pinned `scripts/run\_judge\_eval.py:227--254` instead pools valid leaf labels across examples, and `judge/judge\_eval/evaluate.py:129--134` macro-averages F1 across the binary classes. Equal-paper and pooled-class averaging differ.\par
For example, paper A has 500 positive and 500 negative labels, all correctly classified. Paper B has five of each, all incorrectly classified. Equal-paper F1 is \((1+0)/2=0.5\). After pooling, each class has 500 true positives, five false positives and five false negatives in its one-versus-rest calculation: \[{\rm pooled\ macro\ F1}=\frac{2(500)}{2(500)+5+5}
 =\frac{1000}{1010}=0.990099.\] This illustrates the estimand difference, not actual JudgeEval performance. The code also excludes invalid grades from the classification inputs. Without historical payloads and a publication-time revision, provenance of the original \(0.83\) under the stated paper-level definition is Needs Evidence. Present-source behavior does not prove how the original number was produced.\par
Unweighted F1 also does not by itself bound weighted score error. If \(y_\ell\) are true labels and \(\widehat y_\ell\) are judged labels, then \[|\widehat S-S|
 \le \sum_\ell a_\ell\mathbf1\{\widehat y_\ell\ne y_\ell\}.\] Effective rubric weights matter. JudgeEval provides useful agreement evidence, not a mathematical guarantee of weighted calibration or model-ranking accuracy. The paper discloses nondeterminism, imperfect judging and the need for adversarial testing.\par
A further bounded concern is that the frozen judge prompt in Figures 8–9 describes the provided files as complete and treats missing material as failure, while retrieval may omit files under the context limit. Retrieval omission and true submission absence are different events. Historical error from this mechanism was not measured here.\par
\subsection*{Human baseline and model ordering}
The human comparison is a selected expert baseline: eight ML PhD participants, three initial attempts per paper, continuation of the best after the first week, and a three-paper subset with 48-hour human results versus 36-hour o1 runs. AI assistants are allowed; four attempts used an A100 instead of an A10. These conditions are disclosed on pp. 7–8. The comparison is not a matched-budget effect for randomly sampled humans. No available evidence establishes that matching the conditions would reverse the ordering.\par
The abstract’s unqualified best-tested-agent wording needs the main-setup qualifier already present in the conclusion, p. 9. Table 4 reports \(24.4\%\) for o1 IterativeAgent and \(26.0\%\) with extended runtime, above the Claude BasicAgent headline. Section 5.3 openly explains scaffold-dependent ranking. This is a scope clarification, not concealment of those results.\par
Appendix B filters papers for feasible replication and author collaboration; the corpus is not a random sample of all ML research. Appendix F discloses scaffold prompt development using a small subset of the final dataset. These choices restrict external generalization but do not alone prove corrupted evaluation or false in-sample scores.\par
\section*{Assessment and remaining work}
Object | Judgment | Verified correction or unresolved scope\par
Claude and Gemini means | Minor | Arithmetic contradictions; authoritative historical runs needed.\par
Per-paper SEM | Minor | Sample/population distinction proved; convention should be declared.\par
Aggregate SEM | Needs Evidence | Historical grouping and grade versions unresolved.\par
Log-only result grading | Minor | Current forced-zero path proved; historical effect unknown.\par
JudgeEval aggregation | Needs Evidence | Current statistic differs; original \(0.83\) binding unavailable.\par
Tree score and leaf count | Scoped NIF | Recursion and 8,316 count reconcile.\par
Human comparison | Scope clarification | Disclosed expert selection; no matched effect verified.\par
The supported overall assessment is Minor, with unresolved evidence explicitly separated. No core theoretical guarantee is disproved and no material reversal of the broad empirical conclusions is demonstrated. Reconciliation requires versioned run/grade records, the uncertainty estimator, original JudgeEval aggregation and affected log-only examples. New judge experiments would be needed to quantify score changes. All scientific decisions remain pending, human reviews are empty and author adoption has not been requested.\par
\section*{Calculations, bindings, and production checks}
The decisive `judge/simple.py` has Git blob `1fc0c2b26f5bcc73739b14daa1f2bbd0d36cbe1f`. Checked bytes match their preserved maps. Archives, hierarchy and version identities remain unchanged.\par
Independently authored preserves the full Table 15 transcription, rational means, rounding bounds, SEM variants, leaf counts, the F1 example and empty-list logical condition; results are in . Auxiliary notes preserve additional inspected file hashes, table calculations and locations. These are calculations and static inspections, not empirical reruns.\par
This is an initial review, without acceptance, author contact or publication.\par
9 G. Starace et al., PaperBench: Evaluating AI’s Ability to Replicate AI Research, PMLR 267 (2025). Official publisher PDF, frozen retrieval 25 October 2026. https://proceedings.mlr.press/v267/starace25a.html. OpenAI, frontier-evals, PaperBench source, commit `b78dab1da774a6456b1ed9c48c9ff194ee3ff36d`, 12 October 2026. https://github.com/openai/frontier-evals.\par
\end{document}
